% ===========================================================================
\chapter{Phase 0 --- Fondations et gouvernance de sécurité}
\label{chap:phase0}
% ===========================================================================

\section{Stratégie : la sécurité avant le code}

\subsection{Justification de l'ordonnancement}

La première phase de réalisation ne produit aucune ligne de code applicatif.
Elle met en place la gouvernance de sécurité du dépôt et la pipeline
d'intégration continue, de sorte que \emph{toute} contribution ultérieure soit
contrôlée dès son premier envoi.

Ce choix découle directement de deux principes énoncés au
chapitre~\ref{chap:fondements}.

Le premier est la pratique \pid{P8} \emph{Shift left security}, définie comme
l'intégration des considérations de sécurité le plus tôt possible dans le cycle
de développement. L'étape la plus précoce techniquement atteignable n'est pas
l'intégration continue : c'est le poste du développeur, avant même la création
du commit.

Le second est l'avertissement formulé par l'ASVS~5.0, selon lequel concevoir la
cryptographie au fur et à mesure ou la moderniser après coup coûtera
inévitablement beaucoup plus cher à mettre en œuvre de manière sécurisée que de
l'intégrer dès le départ~\cite{owaspasvs5}. Cette observation se généralise à
l'ensemble des contrôles.

\begin{decision}{pipeline avant code applicatif}
La pipeline d'intégration continue est créée et rendue fonctionnelle
\emph{avant} l'existence du code applicatif.

La conséquence technique est que la pipeline doit réussir sur un projet vide.
Le travail consistant à satisfaire cette contrainte est décrit à la
section~\ref{sec:ci-projet-vide}. Le bénéfice est qu'il devient impossible, par
construction, d'ajouter du code non contrôlé : la règle de protection de branche
dispose d'un nom de contrôle stable dès le premier jour.
\end{decision}

\subsection{Pratiques couvertes par cette phase}

Six pratiques sont mises en œuvre dans cette phase, dont quatre de niveau
Tier~1, sans qu'aucune logique métier n'ait été écrite.

\begin{table}[H]
\centering
\caption{Pratiques couvertes par la phase 0}
\label{tab:phase0-pratiques}
\small
\begin{tabular}{@{}llp{7.4cm}@{}}
\toprule
\textbf{Id.} & \textbf{Niveau} & \textbf{Mise en œuvre} \\
\midrule
\pid{P6}  & Tier 1 & Pipeline GitHub Actions à trois tâches \\
\pid{P8}  & Tier 1 & Crochets \emph{pre-commit} avec Gitleaks \\
\pid{P9}  & Tier 1 & \texttt{CODEOWNERS}, modèle de \emph{pull request} \\
\pid{P11} & Tier 1 & Protection de branche, revue obligatoire \\
\pid{P15} & Tier 2 & Dépôt Git, historique tracé \\
\pid{P30} & Tier 2 & Configuration Gitleaks, règles spécifiques au projet \\
\bottomrule
\end{tabular}
\end{table}

\section{Gouvernance du dépôt}

\subsection{Organisation de l'équipe}

Le dépôt est hébergé sur GitHub à l'adresse
\url{https://github.com/edothecreator/vintage-marketplace-devsecops}.

Les quatre membres du projet y disposent d'un accès en écriture. La répartition
des domaines de responsabilité, formalisée dans le fichier \texttt{CODEOWNERS}
(section~\ref{sec:codeowners}), est la suivante : pipeline et sécurité ;
back-end et API ; front-end et suite de tests d'abus ; infrastructure,
supervision et documentation.

\subsection{Arborescence}

L'arborescence a été définie dès cette phase afin que les artefacts de preuve
disposent d'un emplacement déterminé avant d'être produits.

\begin{lstlisting}[language=bash,caption={Arborescence du dépôt}]
backend/            API REST Express.js + TypeScript
frontend/           front-end web
docs/
  adr/              registres de decisions d'architecture
  threat-model/     modele de menaces et diagrammes
  phases/           suivi des taches par phase
evidence/           preuves par pratique (P1..P11, tier2)
tasks/              fiches de tache pour les membres de l'equipe
.github/workflows/  definitions de la pipeline
\end{lstlisting}

Le répertoire \texttt{evidence/} comporte un sous-répertoire par pratique de
niveau Tier~1. Cette organisation répond à une difficulté pratique identifiée
dès le départ : les journaux d'exécution de pipeline expirent, les branches sont
supprimées, les tableaux de bord sont réinitialisés. Une preuve non capturée au
moment où elle est produite est généralement irrécupérable.

\begin{decision}{capture des preuves intégrée à la définition de « terminé »}
La capture des preuves fait partie de la définition de l'achèvement de chaque
tâche, et non d'une tâche distincte planifiée en fin de projet.
\end{decision}

\section{P6 --- Intégration continue}
\label{sec:ci}

\subsection{Choix de la plate-forme}

L'étude de Prates et Pereira recense cinq outils d'intégration continue :
Jenkins (15 mentions), GitLab CI (8), GitHub Actions (6), Travis CI (3) et
Circle CI (2)~\cite{prates2025}. Plusieurs options étaient donc justifiables au
regard de la source.

\begin{decision}{GitHub Actions (registre ADR 0002)}
GitHub Actions a été retenu pour les motifs suivants :

\begin{itemize}
  \item la plate-forme est native à l'hébergement du dépôt et des
  collaborateurs, ce qui permet l'intégration directe des contrôles obligatoires
  aux \emph{pull requests} --- condition nécessaire pour que \pid{P11} soit
  contraignant et non indicatif ;
  \item aucune infrastructure à maintenir, à l'inverse de Jenkins dont
  l'hébergement, la maintenance et la sécurisation auraient constitué un projet
  en soi, au détriment du temps consacré à l'ingénierie de sécurité ;
  \item l'outil Dependabot est disponible gratuitement, ce qui apporte la
  pratique Tier~2 \pid{P27} \emph{Patch management} à coût quasi nul.
\end{itemize}

GitLab CI, deuxième outil le plus cité, a été écarté car le dépôt et les
collaborateurs sont sur GitHub, et parce que ses fonctionnalités de sécurité les
plus attrayantes relèvent d'offres payantes : nous aurions reconstitué les mêmes
analyses avec les mêmes outils libres, après avoir changé de plate-forme sans
bénéfice.
\end{decision}

\subsection{Structure de la pipeline}

La pipeline de la phase 0 comporte trois tâches.

\begin{table}[H]
\centering
\caption{Tâches de la pipeline en phase 0}
\label{tab:ci-jobs}
\small
\begin{tabular}{@{}llp{6.8cm}@{}}
\toprule
\textbf{Tâche} & \textbf{Pratique} & \textbf{Rôle} \\
\midrule
\texttt{Secret scan (gitleaks)}
& \pid{P30}
& Analyse de l'\emph{historique complet} du dépôt, et non du seul diff \\
\addlinespace
\texttt{Repository hygiene (pre-commit)}
& \pid{P8}
& Exécute en CI les mêmes crochets que localement \\
\addlinespace
\texttt{Build and test}
& \pid{P6}
& Construction et tests ; conditionnelle à l'existence du code \\
\bottomrule
\end{tabular}
\end{table}

Deux choix de configuration méritent d'être explicités.

Le premier concerne les \textbf{permissions du jeton d'exécution}, réduites à la
seule lecture du contenu par défaut. Toute tâche nécessitant davantage doit le
demander explicitement. Il s'agit de l'application du principe de moindre
privilège à la chaîne de production elle-même.

Le second concerne la \textbf{tâche d'hygiène}, qui exécute en intégration
continue exactement les mêmes crochets que ceux installés sur le poste du
développeur. Cette redondance est volontaire : elle garantit que l'option
\texttt{-{}-no-verify}, qui permet localement de contourner un crochet, ne
permet pas de contourner durablement le contrôle, car l'envoi échoue malgré
tout.

\begin{lstlisting}[language=yamlperso,caption={Extrait de
\texttt{.github/workflows/ci.yml} --- principe de moindre privilège et
contrôle de concurrence}]
name: CI

on:
  push:
    branches: ["main", "develop", "feat/**", "fix/**", "demo/**"]
  pull_request:
    branches: ["main", "develop"]

# Moindre privilege par defaut. Toute tache necessitant davantage
# doit le demander explicitement.
permissions:
  contents: read

# Un nouvel envoi sur la meme reference annule l'execution precedente,
# sauf sur main.
concurrency:
  group: ${{ github.workflow }}-${{ github.ref }}
  cancel-in-progress: ${{ github.ref != 'refs/heads/main' }}
\end{lstlisting}

\subsection{Réussir sur un projet vide}
\label{sec:ci-projet-vide}

La contrainte énoncée en début de chapitre --- rendre la pipeline fonctionnelle
avant l'existence du code --- impose que la tâche de construction réussisse sur
un dépôt sans application.

La solution retenue consiste à détecter la présence du code applicatif et à
n'exécuter les étapes Node.js que conditionnellement.

\begin{lstlisting}[language=yamlperso,caption={Détection conditionnelle du code
applicatif}]
- name: Detect application scaffold
  id: detect
  run: |
    if [ -f backend/package.json ]; then
      echo "scaffold=true" >> "$GITHUB_OUTPUT"
    else
      echo "scaffold=false" >> "$GITHUB_OUTPUT"
    fi

# ... etapes Node.js conditionnees par steps.detect.outputs.scaffold

- name: No scaffold yet
  if: steps.detect.outputs.scaffold == 'false'
  run: |
    echo "backend/package.json non present - la phase 2 n'a pas commence."
    echo "Cette tache reussit volontairement afin que la protection de"
    echo "branche dispose d'un controle requis stable des la phase 0."
\end{lstlisting}

\subsection{Résultat}

\begin{verification}{pipeline fonctionnelle}
La pipeline a été exécutée avec succès. Les trois tâches aboutissent, et les
étapes Node.js sont bien ignorées en l'absence de code applicatif, ce qui valide
la logique conditionnelle.
\end{verification}

\captureprete{figures/p06-pipeline-run.png}{Première exécution de la pipeline :
les trois tâches réussissent. Les étapes Node.js sont ignorées, la tâche
« No scaffold yet » confirmant la détection correcte de l'absence de code
applicatif.}

\section{P8 --- Shift left security}
\label{sec:shift-left}

\subsection{Choix du mécanisme de crochets}

\begin{decision}{cadre \emph{pre-commit} (registre ADR 0003)}
Le cadre \texttt{pre-commit}~\cite{precommit} a été retenu plutôt que Husky,
pourtant plus idiomatique dans un projet Node.

Le motif est opérationnel : Gitleaks n'est pas un paquet npm. Avec Husky, il
aurait fallu soit demander à chaque membre de l'équipe d'installer manuellement
le binaire Gitleaks, soit invoquer Docker à chaque commit --- lent, et
inopérant pour quiconque n'a pas Docker en fonctionnement. Le cadre
\texttt{pre-commit} télécharge et épingle lui-même la version du binaire : une
seule commande suffit pour toute l'équipe.
\end{decision}

\subsection{Composition des crochets}

Les crochets installés sont délibérément rapides.

\begin{decision}{rapidité des crochets locaux}
Les analyses lentes --- SAST, SCA, DAST --- restent en intégration continue et
ne sont pas exécutées à chaque commit.

Un crochet lent enseigne aux développeurs l'usage de \texttt{-{}-no-verify}, ce
qui annule le contrôle. La rapidité n'est donc pas ici un critère de confort
mais une condition d'efficacité du contrôle.
\end{decision}

Les crochets retenus sont : Gitleaks (détection de secrets), refus des clés
privées, refus des fichiers volumineux, validation syntaxique YAML et JSON,
normalisation des fins de ligne, et un crochet local refusant explicitement la
mise en index d'un fichier \texttt{.env} réel.

\subsection{Omission volontaire}

Les crochets ESLint, Prettier et de vérification de types n'ont
\emph{pas} été configurés à cette phase.

\begin{decision}{report des crochets de qualité à la phase 2}
Sur un projet vide, ces crochets échoueraient immédiatement. La première
expérience de l'équipe avec l'outillage de sécurité serait alors un chemin de
commit cassé --- le moyen le plus rapide d'enseigner le contournement.

Ils seront ajoutés en phase 2, dès l'existence de \texttt{backend/package.json}.
\end{decision}

\section{P30 --- Gestion des secrets}
\label{sec:secrets}

\subsection{Règles spécifiques au projet}

La configuration Gitleaks part du jeu de règles par défaut maintenu en amont,
puis l'étend de quatre règles propres au projet.

\begin{lstlisting}[caption={Extrait de \texttt{.gitleaks.toml} --- règles
spécifiques au projet}]
[extend]
useDefault = true

[[rules]]
id = "gemini-api-key"
description = "Cle API Google Gemini / AI Studio"
regex = '''AIza[0-9A-Za-z\-_]{35}'''
tags = ["key", "google", "gemini"]

[[rules]]
id = "generic-jwt-secret-assignment"
description = "Secret de signature JWT code en dur"
regex = '''(?i)(jwt[_-]?secret|token[_-]?secret|signing[_-]?key)\s*[:=]\s*['"][^'"\s]{8,}['"]'''
tags = ["secret", "jwt"]

[[rules]]
id = "hardcoded-db-url"
description = "Chaine de connexion avec identifiants en ligne"
regex = '''(?i)(postgres|postgresql|mysql|mongodb(\+srv)?)://[^:\s'"]+:[^@\s'"]+@'''
tags = ["credential", "database"]
\end{lstlisting}

La liste d'autorisation est volontairement restreinte aux valeurs de
substitution, aux fixtures de test et à la documentation. Une liste
d'autorisation trop large produirait une configuration qui semble fonctionner
tout en ne détectant rien.

\subsection{Vérification de l'efficacité des règles}

Une configuration de détection non testée n'est pas un contrôle. Les règles ont
donc été vérifiées par introduction volontaire d'un fichier de test contenant
une fausse clé Gemini et un secret de signature JWT codé en dur.

\begin{verification}{les règles de détection fonctionnent}
Les deux règles ont déclenché. La sortie de l'outil est reproduite ci-dessous,
les valeurs sensibles étant automatiquement masquées par l'option
\texttt{-{}-redact} :

\begin{lstlisting}[basicstyle=\ttfamily\tiny,numbers=none]
gitleaks (block secrets before commit)...................................Failed

Finding:     geminiApiKey: "REDACTED0a",
Secret:      REDACTED
RuleID:      gemini-api-key
Entropy:     5.285402
Tags:        [key google gemini]
File:        backend/src/config/TEMP-leak-test.ts
Line:        3

Finding:     REDACTED: "sup3rS3cr3tSigningKeyForTokens"
Secret:      REDACTED
RuleID:      generic-jwt-secret-assignment
Entropy:     2.725481
Tags:        [secret jwt]
File:        backend/src/config/TEMP-leak-test.ts
\end{lstlisting}

Le fichier de test a été supprimé immédiatement après vérification. Cette
vérification présente un second intérêt : elle valide par avance la
démonstration V2 du chapitre~\ref{chap:vulnerabilites}.
\end{verification}

\section{P9 --- Collaboration}
\label{sec:codeowners}

\subsection{Attribution de la propriété du code}

Le fichier \texttt{CODEOWNERS} associe un ou plusieurs propriétaires à chaque
chemin du dépôt. Combiné à l'option de protection de branche exigeant la revue
des propriétaires, il transforme la revue en exigence structurelle.

Les chemins sensibles à la sécurité --- définitions de pipeline, configuration
Gitleaks, crochets, politique de sécurité --- sont attribués au responsable de
la pipeline, car une modification de ces fichiers peut désactiver silencieusement
un contrôle. Les répertoires d'authentification et d'intergiciels exigent deux
relecteurs.

\begin{constat}{\texttt{CODEOWNERS} échoue silencieusement}
GitHub ignore sans avertissement les noms d'utilisateur inexistants ou ne
disposant pas d'un accès en écriture. Le fichier peut donc avoir toutes les
apparences d'un contrôle fonctionnel sans rien imposer.

Les quatre identifiants ont par conséquent été validés par appel à l'API avant
usage. La validité du fichier est contrôlable par le point d'accès dédié :

\begin{lstlisting}[language=bash,numbers=none]
gh api repos/<owner>/<repo>/codeowners/errors
\end{lstlisting}
\end{constat}

\subsection{Modèle de \emph{pull request}}

Un modèle de \emph{pull request} impose une liste de contrôle de sécurité à
chaque contribution, couvrant la validation des entrées, le contrôle d'accès,
les secrets, les dépendances, la journalisation, et une section spécifique au
code produit avec l'assistance d'une IA générative, conformément à la
section~7 du sujet.

La consigne retenue est qu'une case non cochée doit être justifiée par écrit :
une case laissée vide sans explication est traitée comme non revue.

\captureaplacer{figures/p09-collaborators.png}{Page des collaborateurs du dépôt
(Settings $\rightarrow$ Collaborators) montrant les quatre membres de l'équipe
avec accès en écriture. À capturer après acceptation des invitations.}

\section{P11 --- Revue de code}
\label{sec:revue-code}

\subsection{Protection de la branche principale}

La revue de code n'est une pratique mise en œuvre que si elle est contraignante.
La configuration appliquée à la branche \texttt{main} est la suivante.

\begin{table}[H]
\centering
\caption{Règles de protection appliquées à la branche \texttt{main}}
\label{tab:branch-protection}
\small
\begin{tabular}{@{}p{8.2cm}p{5.4cm}@{}}
\toprule
\textbf{Règle} & \textbf{Valeur} \\
\midrule
Revues approbatrices requises & 1 \\
Revue des propriétaires du code requise & oui \\
Invalidation des revues obsolètes après envoi & oui \\
Approbation du dernier envoi requise & oui \\
Contrôles d'état requis & 3, en mode strict \\
Résolution des discussions requise & oui \\
Envois forcés & interdits \\
Suppression de branche & interdite \\
\bottomrule
\end{tabular}
\end{table}

\begin{constat}{les noms des contrôles requis doivent correspondre exactement}
Les trois contrôles d'état requis sont désignés par le \emph{nom} des tâches de
la pipeline. Un renommage dans \texttt{ci.yml} détache silencieusement la règle
de protection, qui continue d'exister tout en n'exigeant plus rien.

La correspondance exacte entre les noms de tâches effectivement exécutés et les
noms exigés par la protection a donc été vérifiée explicitement.
\end{constat}

\captureaplacer{figures/p11-branch-protection-reviews.png}{Règles de protection
de la branche \texttt{main} (Settings $\rightarrow$ Branches $\rightarrow$ Edit) :
exigence de \emph{pull request}, une revue approbatrice, invalidation des revues
obsolètes, revue des propriétaires du code, approbation du dernier envoi.}

\captureaplacer{figures/p11-branch-protection-status-checks.png}{Suite des règles
de protection : les trois contrôles d'état requis, l'exigence de branche à jour,
et l'interdiction des envois forcés et des suppressions.}

\section{Difficultés rencontrées et décisions induites}

Trois obstacles techniques ont été rencontrés durant cette phase. Ils sont
rapportés ici parce qu'ils ont chacun entraîné une décision documentée, et parce
que leur résolution fait partie du travail d'ingénierie accompli.

\subsection{Protection de branche indisponible sur dépôt privé}

Le dépôt a initialement été créé en mode privé, choix instinctivement plus
prudent. Les tentatives d'activation de la protection de branche, puis des
\emph{rulesets}, ont toutes deux échoué :

\begin{lstlisting}[numbers=none,basicstyle=\ttfamily\scriptsize]
HTTP 403 - Upgrade to GitHub Pro or make this repository public
            to enable this feature.
\end{lstlisting}

L'offre GitHub Free n'autorise ni les branches protégées ni les \emph{rulesets}
sur un dépôt privé. Trois options s'ouvraient : rendre le dépôt public, obtenir
GitHub Pro gratuitement via le \emph{Student Developer Pack}, ou renoncer à la
protection de branche.

La troisième option a été écartée d'emblée : elle aurait réduit la preuve de
deux pratiques de niveau Tier~1 de « la fusion est techniquement bloquée » à
« nous avons rédigé un modèle », ce qui ne démontre pas un contrôle mis en
œuvre.

\begin{decision}{dépôt public (registre ADR 0004)}
Le dépôt a été rendu public. Outre le déblocage de la protection de branche,
ce choix apporte un bénéfice indépendant et matériel : \textbf{les minutes
d'exécution GitHub Actions ne sont pas comptabilisées sur un dépôt public}.

Ce point n'est pas accessoire. Les analyses ZAP complètes et SonarQube sont les
étages les plus coûteux en temps, et le chapitre~\ref{chap:vulnerabilites}
déclenche à lui seul sept exécutions complètes de la pipeline sur sept branches
de démonstration. Sur un dépôt privé, l'allocation gratuite de 2000 minutes
mensuelles nous aurait contraints à rationner les analyses pour des motifs
étrangers à l'ingénierie de sécurité.

La contrepartie est que le code et la configuration sont lisibles publiquement,
ce qui est acceptable pour un travail académique sans contenu propriétaire, et
qu'\emph{aucun} secret réel ne doit jamais entrer dans le dépôt --- ce qui
était déjà la politique, outillée par les crochets locaux et l'analyse
d'historique complet en CI.
\end{decision}

\subsection{Interférence avec la démonstration V2}

Le passage en dépôt public active la protection contre l'envoi de secrets de
GitHub, qui analyse les motifs propres à des fournisseurs identifiés et peut
rejeter l'envoi.

\begin{constat}{un contrôle externe peut masquer le nôtre}
La démonstration V2 prévue au chapitre~\ref{chap:vulnerabilites} introduit
volontairement un secret codé en dur afin que Gitleaks le détecte et fasse
échouer la construction.

Si ce secret est une clé API Google de la forme \texttt{AIza...}, GitHub rejette
l'envoi \emph{avant} que notre pipeline ne s'exécute. Notre contrôle paraîtrait
alors avoir échoué, alors qu'il n'aurait simplement jamais été sollicité.

La démonstration V2 utilise donc le motif de secret de signature JWT, détecté
par notre règle \texttt{generic-jwt-secret-assignment} et non reconnu par
l'analyseur de GitHub. Les deux règles ont été vérifiées
(section~\ref{sec:secrets}).

Cette substitution améliore d'ailleurs la démonstration : un secret de signature
codé en dur est une erreur de développement réaliste, et sa détection repose sur
\emph{notre} configuration plutôt que sur la liste de motifs d'un tiers.
\end{constat}

\subsection{Blocage du gestionnaire d'identifiants Git}

Les opérations Git distantes se sont initialement interrompues sans message.
Le diagnostic a établi que le gestionnaire d'identifiants configuré
(\emph{Git Credential Manager}) tentait d'afficher une boîte de dialogue
d'authentification impossible à présenter dans le contexte d'exécution, et
attendait indéfiniment une saisie.

La résolution a consisté à router les identifiants Git par le jeton de l'outil
en ligne de commande GitHub :

\begin{lstlisting}[language=bash,numbers=none]
gh auth setup-git
\end{lstlisting}

Ce point est consigné car les autres membres de l'équipe rencontreront le même
obstacle au premier clonage. Il figure à ce titre dans la documentation
d'installation.

\section{Registres de décisions d'architecture}

Quatre registres de décisions ont été produits durant cette phase.

\begin{table}[H]
\centering
\caption{Registres de décisions d'architecture de la phase 0}
\label{tab:adr-phase0}
\small
\begin{tabular}{@{}llp{8.2cm}@{}}
\toprule
\textbf{Réf.} & \textbf{Statut} & \textbf{Objet} \\
\midrule
ADR 0001 & Accepté & Tenue de registres de décisions d'architecture \\
ADR 0002 & Accepté & GitHub Actions comme plate-forme CI/CD \\
ADR 0003 & Accepté & Cadre \texttt{pre-commit} pour les crochets \emph{shift left} \\
ADR 0004 & Accepté & Visibilité publique du dépôt et atténuation pour V2 \\
\bottomrule
\end{tabular}
\end{table}

Le registre ADR 0001 justifie la pratique elle-même. Il s'appuie sur
l'observation que le raisonnement conduisant à une décision est ce qui s'efface
le plus vite, et sur l'exigence de l'ASVS~5.0 relative aux décisions de sécurité
documentées. Les registres sont immuables : une décision révisée donne lieu à un
nouveau registre qui remplace le précédent, l'original demeurant afin que
l'historique du raisonnement subsiste.

\section{Bilan de la phase}

\begin{table}[H]
\centering
\caption{État d'avancement à l'issue de la phase 0}
\label{tab:bilan-phase0}
\small
\begin{tabular}{@{}llp{6.6cm}@{}}
\toprule
\textbf{Élément} & \textbf{État} & \textbf{Observation} \\
\midrule
Dépôt, arborescence, configuration & \fait{} & 4 commits \\
Pipeline d'intégration continue    & \fait{} & 2 exécutions, toutes réussies \\
Crochets \emph{pre-commit}         & \fait{} & efficacité vérifiée \\
Règles de détection de secrets     & \fait{} & deux règles déclenchées en test \\
Protection de branche              & \fait{} & vérifiée par API \\
\texttt{CODEOWNERS}                & \encours{} & en attente d'acceptation des invitations \\
Registres de décisions             & \fait{} & ADR 0001 à 0004 \\
Captures de preuve                 & \encours{} & 1 fournie sur 3 \\
\bottomrule
\end{tabular}
\end{table}

L'ingénierie de la phase est achevée. Deux éléments de preuve restent à
collecter : les captures des réglages de plate-forme, qui ne peuvent être
produites que par un administrateur du dépôt et constituent la seule preuve
qu'un contrôle est effectivement activé, et l'acceptation des invitations par
les trois autres membres, préalable au fonctionnement effectif de l'exigence de
revue par les propriétaires du code.

Aucun de ces deux points ne bloque la phase suivante.
